\section{AMC 8 2015}
        \begin{enumerate}
        	\item (\textit{AMC 8 2015}) How many square yards of carpet are required to cover a rectangular floor that is $12$ feet long and $9$ feet wide? (There are $3$ feet in a yard.)


 $\textbf{(A) }12\qquad\textbf{(B) }36\qquad\textbf{(C) }108\qquad\textbf{(D) }324\qquad \textbf{(E) }972$


 
 \textbf{Answer: }A



\textbf{Solution 1}

We need $12\cdot9$ square feet of carpet to cover the floor.  Since there are $9$ square feet in a square yard, we divide this by $9$ to get $\bold{\boxed{\textbf{(A) }12}}$ square yards.

 



\textbf{Solution 2}

Since there are $3$ feet in a yard, we divide $9$ by $3$ to get $3$ yards, and $12$ by $3$ to get $4$ yards. To find the area of the carpet, we then multiply these two values together to get $3\cdot4=\boxed{\textbf{(A) }12}$ square yards.

 


	\item (\textit{AMC 8 2015}) Point $O$ is the center of the regular octagon $ABCDEFGH$, and $X$ is the midpoint of the side $\overline{AB}.$ What fraction of the area of the octagon is shaded?


 
                    \begin{figure}[H]
                        \centering
                        \includegraphics[width=0.5\textwidth]{Problem/AMC_8_AMC_2015_2_Problem_Figure_1.png}
                    \end{figure}
                    


 $\textbf{(A) }\frac{11}{32} \quad\textbf{(B) }\frac{3}{8} \quad\textbf{(C) }\frac{13}{32} \quad\textbf{(D) }\frac{7}{16}\quad \textbf{(E) }\frac{15}{32}$


 
 \textbf{Answer: }D



\textbf{Solution 1}

Since octagon $ABCDEFGH$ is a regular octagon, it is split into $8$ equal parts, such as triangles $\bigtriangleup ABO, \bigtriangleup BCO, \bigtriangleup CDO$, etc.  These parts, since they are all equal, are $\frac{1}{8}$ of the octagon each.  The shaded region consists of $3$ of these equal parts plus half of another, so the fraction of the octagon that is shaded is $\frac{1}{8}+\frac{1}{8}+\frac{1}{8}+\frac{1}{16}=\boxed{\textbf{(D) }\dfrac{7}{16}}.$

 



\textbf{Solution 2}


                    \begin{figure}[H]
                        \centering
                        \includegraphics[width=0.5\textwidth]{Solution/AMC_8_AMC_2015_2_Solution_Figure_1.png}
                    \end{figure}
                    

 The octagon has been divided up into $16$ identical triangles (and thus they each have equal area). Since the shaded region occupies $7$ out of the $16$ total triangles, the answer is $\boxed{\textbf{(D)}~\dfrac{7}{16}}$.

 -Flare

 



\textbf{Solution 3}

For starters, what I find helpful is to divide the whole octagon up into triangles as shown here:
                    \begin{figure}[H]
                        \centering
                        \includegraphics[width=0.5\textwidth]{Solution/AMC_8_AMC_2015_2_Solution_Figure_2.png}
                    \end{figure}
                    

 Now, it is just a matter of counting the larger triangles. Remember that $\triangle BOX$ and $\triangle XOA$ are not full triangles and are only half for these purposes. We count it up and we get a total of $\frac{3.5}{8}$ of the shape shaded. We then simplify it to get our answer of $\boxed{\textbf{(D)}~\frac{7}{16}}$.

 



\textbf{Solution 4}


                    \begin{figure}[H]
                        \centering
                        \includegraphics[width=0.5\textwidth]{Solution/AMC_8_AMC_2015_2_Solution_Figure_2.png}
                    \end{figure}
                    

 We can divide the octagon into 8 parts and pretend that the area is 64. We know that X is the midpoint of BA and that each space between two points is 8 because 64/8=8. This means that BX=4 because 8/2=4. Then, we add that to 3*8 because there are 3 spaces between points that are each 8. After that, you turn it into a fraction, 28/64, and simplify to get  $\boxed{\textbf{(D)}~\frac{7}{16}}$.

 



\textbf{Solution 5 (Quick and Efficient)}


                    \begin{figure}[H]
                        \centering
                        \includegraphics[width=0.5\textwidth]{Solution/AMC_8_AMC_2015_2_Solution_Figure_4.png}
                    \end{figure}
                    We see that the midpoint of $\overline{AB}$ is $X$. We can also notice that point $X$ to point $E$ *ALMOST* makes the octagon half shaded. Since we know an octagon has 8 sides, $X$ is one-half of $\overline{AB}$ which means instead of counting in eights, we need to count in sixteenths. Now back to the problem. Let's say that the midpoint of $\overline{EF}$ is $Y$, then $\overline{XY}$ cuts the octagon perfectly in half, but since there is no $Y$, we are now $\frac1{16}$ away from one-half. We know that $\frac1{2}=\frac8{16}$, so we can just do $\frac8{16}-\frac1{16}$ which gets us, for the final answer,  $\boxed{\textbf{(D)}~\frac{7}{16}}$.

 


	\item (\textit{AMC 8 2015}) Jack and Jill are going swimming at a pool that is one mile from their house. They leave home simultaneously. Jill rides her bicycle to the pool at a constant speed of $10$ miles per hour. Jack walks to the pool at a constant speed of $4$ miles per hour. How many minutes before Jack arrives does Jill arrive?


 $\textbf{(A) }5\qquad\textbf{(B) }6\qquad\textbf{(C) }8\qquad\textbf{(D) }9\qquad \textbf{(E) }10$


 
 \textbf{Answer: }D



\textbf{Solution}

Using $d=rt$, we can set up an equation for when Jill arrives at the swimming pool:

 $1=10t$

 Solving for $t$, we get that Jill gets to the pool in $\frac{1}{10}$ of an hour, which is $6$ minutes.  Doing the same for Jack, we get that 

 Jack arrives at the pool in $\frac{1}{4}$ of an hour, which in turn is $15$ minutes.  Thus, Jill has to wait $15-6=\boxed{\textbf{(D)}~9}$ 

 minutes for Jack to arrive at the pool.

 



\textbf{Solution 2}

T=D/sJill: (1/10)x60 because in minutes, is equal to 6 minJack:(1/4)x60 is 15 minutes.15-6 is 9, so our answer is $15-6=\boxed{\textbf{(D)}~9}$  - TheNerdWhoIsNerdy.

 


	\item (\textit{AMC 8 2015}) The Centerville Middle School chess team consists of two boys and three girls. A photographer wants to take a picture of the team to appear in the local newspaper. She decides to have them sit in a row with a boy at each end and the three girls in the middle. How many such arrangements are possible? 


 $\textbf{(A) }2\qquad\textbf{(B) }4\qquad\textbf{(C) }5\qquad\textbf{(D) }6\qquad \textbf{(E) }12$


 
 \textbf{Answer: }E



\textbf{Solution}

There are $2! = 2$ ways to order the boys on the ends, and there are $3!=6$ ways to order the girls in the middle. We get the answer to be $2 \cdot 6 = \boxed{\textbf{(E) }12}$.

 


	\item (\textit{AMC 8 2015}) Billy's basketball team scored the following points over the course of the first $11$ games of the season: 
 \[42, 47, 53, 53, 58, 58, 58, 61, 64, 65, 73\] If his team scores 40 in the 12th game, which of the following statistics will show an increase?


 $\textbf{(A) } \text{range} \qquad \textbf{(B) } \text{median} \qquad \textbf{(C) } \text{mean} \qquad \textbf{(D) } \text{mode} \qquad \textbf{(E) } \text{mid-range}$


 
 \textbf{Answer: }A



\textbf{Solution 1}

When they score a $40$ on the next game, the range increases from $73-42=31$ to $73-40=33$.  This means the $\boxed{\textbf{(A) } \text{range}}$ increased.

 Note: The range is defined to be the difference of the largest element and the smallest element of a set.

 



\textbf{Solution 2}

Because $40$ is less than the score of every game they've played so far, the measures of center will never rise. Only measures of spread, such as the $\boxed{\textbf{(A)}~\text{range}}$, may increase.

 


	\item (\textit{AMC 8 2015}) In $\triangle ABC$, $AB=BC=29$, and $AC=42$. What is the area of $\triangle ABC$?


 $\textbf{(A) }100\qquad\textbf{(B) }420\qquad\textbf{(C) }500\qquad\textbf{(D) }609\qquad \textbf{(E) }701$


 
 \textbf{Answer: }B



\textbf{Solution 1}

We know the semi-perimeter of $\triangle ABC$ is $\frac{29+29+42}{2}=50$. Next, we use Heron's Formula to find that the area of the triangle is just $\sqrt{50(50-29)^2(50-42)}=\sqrt{50 \cdot 21^2 \cdot 8}=\boxed{\textbf{(B) }420}$.

 



\textbf{Solution 2 (easier)}

Splitting the isosceles triangle in half, we get a right triangle with hypotenuse $29$ and leg $21$. Using the Pythagorean Theorem , we know the height is $\sqrt{29^2-21^2}=20$. Now that we know the height, the area is$\dfrac{(20)(42)}{2} = \boxed{\textbf{(B)}~420}$.

 


	\item (\textit{AMC 8 2015}) Each of two boxes contains three chips numbered $1$, $2$, $3$. A chip is drawn randomly from each box and the numbers on the two chips are multiplied. What is the probability that their product is even?


 $\textbf{(A) }\frac{1}{9}\qquad\textbf{(B) }\frac{2}{9}\qquad\textbf{(C) }\frac{4}{9}\qquad\textbf{(D) }\frac{1}{2}\qquad \textbf{(E) }\frac{5}{9}$


 
 \textbf{Answer: }E



\textbf{Solution 1}

(This solution is similar to Solution 2.)Let's make this a problem with boxes. 

 In total, there are 9 products derived from these numbers (because 3 numbers per box). This would be our denominator. 

 Every time a spinner lands on 2, we get an even product. These are $(1,2)$, $(2,1)$, $(2,2)$, $(2,3)$, and finally $(3,2)$. 

 Going back, we see there are $3\cdot3=9$ possible combinations, and we have $5$ evens, the final answer is  

 $\boxed{\textbf{(E) }\frac{5}{9}}$.

 



\textbf{Solution 2}

You can also make this problem into a spinner problem. You have the first spinner with $3$ equally divided 

 sections: $1, 2$, and $3$. You make a second spinner that is identical to the first, with $3$ equal sections of 

 $1$,$2$, and $3$. If the first spinner lands on $1$, it must land on two for the result to be even. You write down the first 

 combination of numbers: $(1,2)$. Next, if the spinner lands on $2$, it can land on any number on the second 

 spinner. We now have the combinations of $(1,2) ,(2,1), (2,2),$ and $(2,3)$. Finally, if the first spinner ends on $3$, we 

 have $(3,2).$ Since there are $3\cdot3=9$ possible combinations, and we have $5$ evens, the final answer is  

 $\boxed{\textbf{(E) }\frac{5}{9}}$.

 



\textbf{Solution 3}

We can also list out the numbers. Box A has chips $1$, $2$, and $3$, and Box B also has chips $1$, $2$, and $3$. Chip $1$ (from Box A) 

 could be with 3 partners from Box B. This is also the same for chips $2$ and $3$ from Box A. $3+3+3=9$ total sums. Chip $1$ could be multiplied with 2 other chips to make an even product, just like chip $3$. Chip $2$ can only multiply with 1 chip.  $2+2+1=5$. The answer is $\boxed{\textbf{(E) }\frac{5}{9}}$.

 



\textbf{Solution 4}

Here is another way:

 Let's start by finding the denominator: Total choices.There are $3$ chips we can choose from in the 1st box, and $3$ chips we can choose from in the 2nd box. We do $3*3$, and get $9$.Now - to find the numerator: Desired choices.To get an even number, we need to pick 2 from at least one of the boxes. There are $2$ choices as to finding which box we will draw the 2 from. Then, we have $3$ choices from the other box to pick any of the other chips, $1, 2,$ and $3$. 

 $\frac{3 \cdot2}{9} = \frac{6}{9}$

 However, we are over counting the $(2,2)$ configuration twice, and so, we subtract that one configuration from our total.$\frac{6}{9} - \frac{1}{9}$.Thus, our answer is $\boxed{\textbf{(E) }\frac{5}{9}}$.

 



\textbf{Solution 5}

This might take longer to solve, but you definitely will have the right answer.You first list out all the possible combinations regardless if the product is even or odd.$1 \cdot 1 = 1.$ $1 \cdot 2 = 2.$ $1 \cdot 3 = 3.$ $2 \cdot 1 = 2.$ $2 \cdot 2 = 4.$ $2 \cdot 3 = 6.$ $3 \cdot 1 = 3.$ $3 \cdot 2 = 6.$ $3 \cdot 3 = 9.$   There are 9 possible combinations and $5$ of the combinations’ products are even.  So, our answer is $\boxed{\textbf{(E) }\frac{5}{9}}$.

 



\textbf{Solution 6}

We can use complementary counting to solve this problem. If we want to find the number of ways for the product to be odd, then the chip chosen from either box has to be $1$ or $3$. Since there are 2 ways out of 3 to pick an even chip from a box, the probability of picking an even chip from a box is $\frac{2}{3}$. Since there are $2$ boxes, we can find that the answer is $1-\frac{2}{3}\cdot\frac{2}{3}$=$\boxed{\textbf{(E) }\frac{5}{9}}$.

 


	\item (\textit{AMC 8 2015}) What is the smallest whole number larger than the perimeter of any triangle with a side of length $5$ and a side of length $19$?


 $\textbf{(A) }24\qquad\textbf{(B) }29\qquad\textbf{(C) }43\qquad\textbf{(D) }48\qquad \textbf{(E) }57$


 
 \textbf{Answer: }D



\textbf{Solution}

We know from the \href{https://artofproblemsolving.com/wiki/index.php?title=Triangle\_Inequality}{Triangle Inequality} that the last side, $s$, fulfills $s<5+19=24$. Adding $5+19$ to both sides of the inequality, we get $s+5+19<48$, and because $s+5+19$ is the perimeter of our triangle, $\boxed{\textbf{(D)}\ 48}$ is our answer.

 


 



\textbf{Solution 2}

$s$, fulfills $s<5+19=24$, meaning that s<24. The greatest value of s can be 23. 23+5+19 is 47. Since 48 is our nearest answer choice, the answer is $\boxed{\textbf{(D)}\ 48}$.-TheNerdWhoIsNerdy.

 minor edit by themathkidwoohooooo:the greatest value of s is 24-$e$, where $e$ represents an infinitesimal (a constant smaller than any positive number but not zero or negative. it is between the positives and 0. another way to think of it is $\frac {1} {\infty}$). NOT 23!!!!!!!!!!!!!!!

 


	\item (\textit{AMC 8 2015}) On her first day of work, Janabel sold one widget. On day two, she sold three widgets. On day three, she sold five widgets, and on each succeeding day, she sold two more widgets than she had sold on the previous day. How many widgets in total had Janabel sold after working $20$ days?


 $\textbf{(A) }39\qquad\textbf{(B) }40\qquad\textbf{(C) }210\qquad\textbf{(D) }400\qquad \textbf{(E) }401$


 
 \textbf{Answer: }D



\textbf{Solution 1}

First, we recognize that the number of widgets Janabel sells each day forms an arithmetic sequence. On day 1, she sells 1 widget, on day 2, she sells 3 widgets, on day 3, she sells 5 widgets, and so on. This sequence can be described by the first term \( a = 1 \) and the common difference \( d = 2 \).

 Step 1: Find the number of widgets sold on day 20

 The formula for the \( n \)-th term of an arithmetic sequence is:
 \[a_n = a + (n-1)d\]For \( n = 20 \):
 \[a_{20} = 1 + (20-1) \cdot 2 = 1 + 38 = 39\]

 So, Janabel sells 39 widgets on day 20.

 Step 2: Calculate the total number of widgets sold in 20 days}

 The sum \( S_n \) of the first \( n \) terms of an arithmetic sequence is given by:
 \[S_n = \frac{n}{2} \cdot (a + a_n)\]

 For \( n = 20 \):
 \[S_{20} = \frac{20}{2} \cdot (1 + 39) = 10 \cdot 40 = 400\]

 Therefore, the total number of widgets sold after 20 days is:$(40\cdot10)=\boxed{\textbf{(D)}~400}$

 



\textbf{Solution 2}

First, we have to find how many widgets she makes on Day $20$. We can write the linear equation $y=-1+2x$ to represent this situation. Then, we can plug in $20$ for $x$:  $y=-1+2(20)$ -- $y=-1+40$ -- $y=39$. The sum of $1,3,5, ... 39$ is $\dfrac{(1 + 39)(20)}{2}= \boxed{\textbf{(D)}~400}$.

 



\textbf{Solution 3 ( !!FASTEST SOLUTION!! )}

The sum is just the first $20$ odd counting/natural numbers, which is $20^2=\boxed{\textbf{(D)}~400}$.

 Note: The sum of the first $x$ odd numbers is $x^2$.

 



\textbf{Solution 4}

We can easily find out she makes $2\cdot20-1 = 39$ widgets on Day $20$. Then, we make the sum of $1,3, 5, ... ,35,37,39$ by adding in this way: $(1+39)+(3+37)+(5+35)+...+(19+21)$, which include $10$ pairs of $40$. So, the sum of $1,3,5, ...~39$ is $(40\cdot10)=\boxed{\textbf{(D)}~400}$.--LarryFlora

 


	\item (\textit{AMC 8 2015}) How many integers between $1000$ and $9999$ have four distinct digits?


 $\textbf{(A) }3024\qquad\textbf{(B) }4536\qquad\textbf{(C) }5040\qquad\textbf{(D) }6480\qquad \textbf{(E) }6561$


 
 \textbf{Answer: }B



\textbf{Solution 1}

There are $9$ choices for the first number, since it cannot be $0$, there are only $9$ choices left for the second number since it must differ from the first, $8$ choices for the third number, since it must differ from the first two, and $7$ choices for the fourth number, since it must differ from all three.  This means there are $9 \times 9 \times 8 \times 7=\boxed{\textbf{(B) }4536}$ integers between $1000$ and $9999$ with four distinct digits.

 
If you're too lazy to do $9 \times 9 \times 8 \times 7= 4536$, we can notice $9 \times 9 = 81$ and $8 \times 7=56$ The unit digit must be 6. Therefore, the answer is the only answer choice with a unit digit of 6, which is $\boxed{B}$.

 


	\item (\textit{AMC 8 2015}) In the small country of Mathland, all automobile license plates have four symbols. The first symbol must be a vowel ($A, E, I, O,$ or $U$), the second and third symbols must be two different letters among the $21$ non-vowels in the alphabet, and the fourth symbol must be a digit ($0$ through $9$). If the symbols are chosen at random subject to these conditions, what is the probability (In fractions) that the plate will read "$AMC8$"?


 $\textbf{(A) } \frac{1}{22,050} \qquad \textbf{(B) } \frac{1}{21,000}\qquad \textbf{(C) } \frac{1}{10,500}\qquad \textbf{(D) } \frac{1}{2,100} \qquad \textbf{(E) } \frac{1}{1,050}$


 
 \textbf{Answer: }B



\textbf{Solution 1}

There is one favorable case, which is the license plate says "AMC8."  We must now find how many total cases there are.  There are $5$ choices for the first letter (since it must be a vowel), $21$ choices for the second letter (since it must be of $21$ consonants), $20$ choices for the third letter (since it must differ from the second letter), and $10$ choices for the digit.  This leads to $5 \cdot 21 \cdot 20 \cdot 10=21,000$ total possible license plates.  Therefore, the probability of a license plate saying "AMC8" is $\boxed{\textbf{(B) } \frac{1}{21,000}}$.

 



\textbf{Solution 2}

The probability of choosing A as the first letter is $\dfrac{1}{5}$.  The probability of choosing $M$ next is $\dfrac{1}{21}$.  The probability of choosing C as the third letter is $\dfrac{1}{20}$ (since there are $20$ other consonants to choose from other than M). The probability of having $8$ as the last number is $\dfrac{1}{10}$.  We multiply all these to obtain $\dfrac{1}{5}\cdot \dfrac{1}{21} \cdot \dfrac{1}{20}\cdot \dfrac{1}{10}=\dfrac{1}{5\times 21\times 20\times 10}=\dfrac{1}{21\times 100\times 10}=\boxed{\textbf{(B)}~\dfrac{1}{21,000}}$.

 


	\item (\textit{AMC 8 2015}) How many pairs of parallel edges, such as $\overline{AB},$ and $\overline{GH},$ or $\overline{EH},$ and $\overline{FG},$ does a cube have? Repeating parallel edges do not count toward the final answer.


 
                    \begin{figure}[H]
                        \centering
                        \includegraphics[width=0.5\textwidth]{Problem/AMC_8_AMC_2015_12_Problem_Figure_1.png}
                    \end{figure}
                    


 $\textbf{(A) }6 \qquad\textbf{(B) }12 \qquad\textbf{(C) } 18 \qquad\textbf{(D) } 24 \qquad \textbf{(E) } 36$


 
 \textbf{Answer: }C



\textbf{Solution 1}

We first count the number of pairs of parallel lines that are in the same direction as $\overline{AB}$.  The pairs of parallel lines are $\overline{AB}\text{ and }\overline{EF}$, $\overline{CD}\text{ and }\overline{GH}$, $\overline{AB}\text{ and }\overline{CD}$, $\overline{EF}\text{ and }\overline{GH}$, $\overline{AB}\text{ and }\overline{GH}$, and $\overline{CD}\text{ and }\overline{EF}$.  These are $6$ pairs total.  We can do the same for the lines in the same direction as $\overline{AE}$ and $\overline{AD}$.  This means there are $6\cdot 3=\boxed{\textbf{(C) } 18}$ total pairs of parallel lines.

 



\textbf{Solution 2}

Look at any edge, let's say $\overline{AB}$. There are three ways we can pair $\overline{AB}$ with another edge. $\overline{AB}\text{ and }\overline{EF}$, $\overline{AB}\text{ and }\overline{HG}$, and $\overline{AB}\text{ and }\overline{DC}$. There are 12 edges on a cube. 3 times 12 is 36. We have to divide by 2 because every pair is counted twice, so $\frac{36}{2}$ is $\boxed{\textbf{(C) } 18}$ total pairs of parallel lines.

 



\textbf{Solution 3}

We can use the feature of 3-Dimension in a cube to solve the problem systematically. For example, in the 3-D of the cube, $\overline{AB}$, $\overline{BC}$, and $\overline{BF}$ have $4$ different parallel edges respectively. So it gives us the total pairs of parallel lines are $\binom{4}{2}\cdot3 =\boxed{\textbf{(C) } 18}$.

 --LarryFlora

 



\textbf{Solution 4}

Our first case are lines on the same plane. There are two pairs of parallel lines for a square, and there are 6 squares, which give us $6\cdot2=12$ cases.

 The next cases we need to analyze are the parallel lines on opposite planes (lines such as $\overline{EF}\text{ and } \overline{DC}$). There are 2 pairs for every 2 planes, and there are 6 planes in total, giving us $2\cdot3=6$ cases.

 Adding up the cases, we get $12+6 = \boxed{\textbf{(C) } 18}$

 


	\item (\textit{AMC 8 2015}) How many subsets of two elements can be removed from the set $\{1, 2, 3, 4, 5, 6, 7, 8, 9, 10, 11\}$ so that the mean (average) of the remaining numbers is $6$?


 $\textbf{(A)}\text{ 1}\qquad\textbf{(B)}\text{ 2}\qquad\textbf{(C)}\text{ 3}\qquad\textbf{(D)}\text{ 5}\qquad\textbf{(E)}\text{ 6}$


 
 \textbf{Answer: }D



\textbf{Solution 1}

Since there will be $9$ elements after removal, and their mean is $6$, we know their sum is $54$. We also know that the sum of the set pre-removal is $66$. Thus, the sum of the $2$ elements removed is $66-54=12$. There are only $\boxed{\textbf{(D)}~5}$ subsets of $2$ elements that sum to $12$: $\{1,11\}, \{2,10\}, \{3, 9\}, \{4, 8\}, \{5, 7\}$.

 



\textbf{Solution 2}

We can simply remove $5$ subsets of $2$ numbers while leaving only $6$ behind. The average of this one-number set is still $6$, so the answer is $\boxed{\textbf{(D)}~5}$.

 -tryanotherangle

 


	\item (\textit{AMC 8 2015}) Which of the following integers cannot be written as the sum of four consecutive odd integers?


 $\textbf{(A)}\text{ 16}\qquad\textbf{(B)}\text{ 40}\qquad\textbf{(C)}\text{ 72}\qquad\textbf{(D)}\text{ 100}\qquad\textbf{(E)}\text{ 200}$


 
 \textbf{Answer: }D



\textbf{Solution 1}

Let our $4$ numbers be $n, n+2, n+4, n+6$, where $n$ is odd. Then, our sum is $4n+12$. The only answer choice that cannot be written as $4n+12$, where $n$ is odd, is $\boxed{\textbf{(D)}\text{ 100}}$.

 



\textbf{Solution 2}

If the four consecutive odd integers are $2n-3,~ 2n-1, ~2n+1$ and $2n+3$; then, the sum is $8n$. All the integers are divisible by $8$ except $\boxed{\textbf{(D)}~100}$.

 



\textbf{Solution 3}

If the four consecutive odd integers are $a,~ a+2, ~a+4$ and $a+6$, the sum is $4a+12$, and $4a+12$ divided by $4$ gives $a+3$. This means that $a+3$ must be even. The only integer that does not give an even integer when divided by $4$ is $100$, so the answer is $\boxed{\textbf{(D)}~100}$.

 



\textbf{Solution 4}

From Solution 1, we have the sum of the $4$ numbers to be equal to $4n + 12$. Taking mod 8 gives us $4n + 4 \equiv b \pmod8$ for some residue $b$ and for some odd integer $n$. Since $n \equiv 1 \pmod{2}$, we can express it as the equation $n = 2a + 1$ for some integer $a$. Multiplying 4 to each side of the equation yields $4n = 8a + 4$, and taking mod 8 gets us $4n \equiv 4 \pmod{8}$, so $b = 0$. All the answer choices except choice D is a multiple of 8, and since 100 satisfies all the conditions the answer is $\boxed{\textbf{(D)}~100}$.

 



\textbf{Solution 5}

Since they want CONSECUTIVE odd numbers, it won't be hard to just list the sums out:
 \[16=1+3+5+7\]
 \[40=7+9+11+13\]
 \[72=15+17+19+21\]
 \[200=47+49+51+53\]All of the answer choices can be a sum of consecutive odd numbers except $100$, so the answer is $\boxed{\textbf{(D)} 100}$.

 


	\item (\textit{AMC 8 2015}) At Euler Middle School, $198$ students voted on two issues in a school referendum with the following results: $149$ voted in favor of the first issue and $119$ voted in favor of the second issue. If there were exactly $29$ students who voted against both issues, how many students voted in favor of both issues?


 $\textbf{(A) }49\qquad\textbf{(B) }70\qquad\textbf{(C) }79\qquad\textbf{(D) }99\qquad \textbf{(E) }149$


 
 \textbf{Answer: }D



\textbf{Solution 1}

Let:

 
                \begin{itemize}
                    \item $A$ be the number of students who voted in favor of the first issue,\item $B$ be the number of students who voted in favor of the second issue,\item $A \cap B$ be the number of students who voted in favor of both issues.
                \end{itemize}


                We are given:

 
                \begin{itemize}
                    \item $A = 149$\item $B = 119$\item 29 students voted against both issues, so the number of students who voted for at least one issue is:
                \end{itemize}


                
 \[198 - 29 = 169\]

 By the principle of inclusion and exclusion:
 \[A \cup B = A + B - A \cap B\]Substitute known values:
 \[169 = 149 + 119 - A \cap B\]
 \[169 = 268 - A \cap B\]
 \[A \cap B = 268 - 169 = \boxed{\textbf{(D) }99}\]

 



\textbf{Solution 2}

We can see that this is a Venn Diagram Problem.

 First, we analyze the information given. There are $198$ students. Let's use A as the first issue and B as the second issue. 

 $149$ students were for A, and $119$ students were for B. There were also $29$ students against both A and B. 

 Solving this without a Venn Diagram, we subtract $29$ away from the total, $198$. Out of the remaining $169$ , we have $149$ people for A and $119$ people for B. We add this up to get $268$ . Since that is more than what we need, we subtract $169$ from $268$ to get 

 $\boxed{\textbf{(D)}~99}$.

 
                    \begin{figure}[H]
                        \centering
                        \includegraphics[width=0.5\textwidth]{Solution/AMC_8_AMC_2015_15_Solution_Figure_1.png}
                    \end{figure}
                    

 Note: One could use the Principle of Inclusion-Exclusion in a similar way to achieve the same result.

 


 



\textbf{Solution 3}

There are $198-149=49$ students who voted against the first issue. Out of these, $49-29=20$ voted for the second issue. We also find that $50$ students voted against the second issue but in favor of the first in a similar way. This means that $149-50=119-20=\boxed{\textbf{(D)}~99}$ students voted in favor of both issues.

 


 


	\item (\textit{AMC 8 2015}) In a middle-school mentoring program, a number of the sixth graders are paired with a ninth-grade student as a buddy. No ninth grader is assigned more than one sixth-grade buddy. If $\frac{1}{3}$ of all the ninth graders are paired with $\frac{2}{5}$ of all the sixth graders, what fraction of the total number of sixth and ninth graders have a buddy?


 $\textbf{(A) } \frac{2}{15} \qquad \textbf{(B) } \frac{4}{11} \qquad \textbf{(C) } \frac{11}{30} \qquad \textbf{(D) } \frac{3}{8} \qquad \textbf{(E) } \frac{11}{15}$


 
 \textbf{Answer: }B



\textbf{Solution 1 (Using equations)}

Let the number of sixth graders be $s$, and the number of ninth graders be $n$. Thus, $\frac{n}{3}=\frac{2s}{5}$, which simplifies to $n=\frac{6s}{5}$. Since we are trying to find the value of $\frac{\frac{n}{3}+\frac{2s}{5}}{n+s}$, we can just substitute $\frac{6s}{5}$ for  $n$  into the equation. We then get a value of $\frac{\frac{\frac{6s}{5}}{3}+\frac{2s}{5}}{\frac{6s}{5}+s} = \frac{\frac{6s+6s}{15}}{\frac{11s}{5}} = \frac{\frac{4s}{5}}{\frac{11s}{5}} = \boxed{\textbf{(B)}~\frac{4}{11}}$.

 



\textbf{Solution 2 (Easy)}

We see that the minimum number of ninth graders is $6$, because if there are $3$ then there is $1$ ninth-grader with a buddy, which would mean there are $2.5$ sixth graders, which is impossible (of course unless you really do have half of a person). With $6$ ninth-graders, $2$ of them are in the buddy program, so there $\frac{2}{\frac{2}{5}}=5$ sixth-graders total, two of whom have a buddy. Thus, the desired fraction is $\frac{2+2}{5+6}=\boxed{\textbf{(B) }\frac{4}{11}}$.

 


 


	\item (\textit{AMC 8 2015}) Jeremy's father drives his son to school in rush hour traffic in $20$ minutes. One day, there is no traffic, so his father can drive him $18$ miles per hour faster and gets him to school in $12$ minutes. How far in miles is it to school?


 $\textbf{(A) } 4 \qquad \textbf{(B) } 6 \qquad \textbf{(C) } 8 \qquad \textbf{(D) } 9 \qquad \textbf{(E) } 12$


 
 \textbf{Answer: }D



\textbf{Solution 1}

For starters, we identify d as distance and v as velocity (speed)

 Writing the equation gives us: $\frac{d}{v}=\frac{1}{3}$ and $\frac{d}{v+18}=\frac{1}{5}$.

 This gives $d=\frac{1}{5}v+3.6=\frac{1}{3}v$, which gives $v=27$, which then gives $d=\boxed{\textbf{(D)}~9}$.

 



\textbf{Solution 2}

$d = rt$, $d=\frac{1}{3} \times r = \frac{1}{5} \times (r + 18)$

 $\frac{r}{3} = \frac{r}{5} + \frac{18}{5}$

 
$10r = 270$ so $r = 27$, plug into the first one and it's $\boxed{\textbf{(D)}~9}$ miles to school.

 



\textbf{Solution 3}

We set up an equation in terms of $d$ the distance and $x$ the speed In miles per hour. We have $d=\left (\dfrac{1}{3}\right )(x)=\left (\dfrac{1}{5}\right) (x+18)$, giving
 \[(5)(x)=(3)(x+18)\]
 \[5x=3x+54\]
 \[2x=54\]
 \[x=27\]Hence, $d=\dfrac{27}{3}=\boxed{\textbf{(D)}~9}$.

 



\textbf{Solution 4}

Since it takes $\frac{3}{5}$ of the original time for him to get to school when there is no traffic, the speed must be $\frac{5}{3}$ of the speed in no traffic or $\frac{2}{3}$ more. Letting $x$ be the rate and we know that $\frac{5}{3}x = x + 18$, so we have $\frac{2x}{3} = 18$ miles per hour. Solving for $x$ gives us $27$ miles per hour. Because $20$ minutes is a third of an hour, the distance would then be $9$ miles $\boxed{\textbf{(D)}~9}$.

 



\textbf{Solution 5}

When driving in rush hour traffic, he drives $20$ minutes for one distance ($1d$) to the school. It means he drives $60$ minutes for $3$ distances ($3d$) to the school. When driving in no traffic hours, he drives $12$ minutes for one distance ($1d$) to the school. It means he drives $60$ minutes for $5$ distances ($5d$) to the school. Subtracting these two situations, it gives us $5d-3d = 18 = 2d$, then $d=\frac{18}{2}=9$. So the distance to the school would be $\boxed{\textbf{(D)}~9}$ miles. ~LarryFlora

 



\textbf{Solution 6}

(Ratios)$\textbf{requires edits}$

 In rush hour traffic, we can create this ratio for each distance we are going to try supposing the distance is $d$:

 $d$ : $20$ hours

 In no traffic, we can do the same:

 $d$ : $12$ minutes

 We want the ratio to be the distance to $60$ minutes:

 $d$ : $20$ minutes = $3$d : $60$ minutes 

 $d$ : $12$ minutes = $5$d : $60$ minutes

 This shows the speed in miles per hour. This means that the speed in the 2nd ratio has to be 18 faster than the first ratio. We get

 $3d+18 = 5d \cdot d=9$, or $\boxed{\textbf{(D)}~9}$.

 



\textbf{Solution 7 (Solution 3 but explained better)}

Distance equals speed multiplied by time. Let x/60 represent the speed when Jeremy's father is driving through traffic. Let (x+18)/60 represent the speed when there is no traffic. Because distance equals speed multiplied by time, 20x/60 would represent the distance to school. The "20" in this expression represents the 20 minutes it takes Jeremy's father to drive to school through the traffic, and "60" means 60 minutes. Anyway, 20x/60 simplifies to x/3. Using the same distance equation, the expression 12(x+18)/60 also represents the distance to school. (x+18)/60 represents the speed when there is no traffic, and the "12" means 12 minutes it takes to get to school when there is no traffic. 12(x+18)/60 simplifies to (x+18)/5. Now, since both expressions equate to the distance to school, you get the equation x/3=(x+18)/5. Solving for x gives you 27. This means that Jeremy's father drives 27 miles per hour in traffic. To find the distance to school, simply take 27 and plug it into the expressions (x+18)/5 or x/3. Both give you 9, which is the distance to school.~Evžen D.

 


	\item (\textit{AMC 8 2015}) An arithmetic sequence is a sequence in which each term after the first is obtained by adding a constant to the previous term. For example, $2,5,8,11,14$ is an arithmetic sequence with five terms, in which the first term is $2$ and the constant added is $3$. Each row and each column in this $5\times5$ array is an arithmetic sequence with five terms. The square in the center is labelled $X$ as shown. What is the value of $X$?


 
                    \begin{figure}[H]
                        \centering
                        \includegraphics[width=0.5\textwidth]{Problem/AMC_8_AMC_2015_18_Problem_Figure_1.png}
                    \end{figure}
                    


 $\textbf{(A) }21\qquad\textbf{(B) }31\qquad\textbf{(C) }36\qquad\textbf{(D) }40\qquad \textbf{(E) }42$


 
 \textbf{Answer: }B



\textbf{Solution 1}

We begin filling in the table.  The top row has a first term $1$ and a fifth term $25$, so we have the common difference is $\frac{25-1}4=6$.  This means we can fill in the first row of the table:
                    \begin{figure}[H]
                        \centering
                        \includegraphics[width=0.5\textwidth]{Solution/AMC_8_AMC_2015_18_Solution_Figure_1.png}
                    \end{figure}
                    

 The fifth row has a first term of $17$ and a fifth term of $81$, so the common difference is $\frac{81-17}4=16$.  We can fill in the fifth row of the table as shown:
                    \begin{figure}[H]
                        \centering
                        \includegraphics[width=0.5\textwidth]{Solution/AMC_8_AMC_2015_18_Solution_Figure_2.png}
                    \end{figure}
                    

 We must find the third term of the arithmetic sequence with a first term of $13$ and a fifth term of $49$.  The common difference of this sequence is $\frac{49-13}4=9$, so  the third term is $13+2\cdot 9=\boxed{\textbf{(B) }31}$.

 



\textbf{Solution 2}

The middle term of the first row is $\frac{25+1}{2}=13$, since the middle number is just the average in an arithmetic sequence. Similarly, the middle of the bottom row is $\frac{17+81}{2}=49$. Applying this again for the middle column, the answer is $\frac{49+13}{2}=\boxed{\textbf{(B)}~31}$.

 



\textbf{Solution 3}

The value of $X$ is simply the average of the average values of both diagonals that contain $X$. This is $\frac{\frac{1+81}{2}+\frac{17+25}{2}}{2} =\frac{\frac{82}{2}+\frac{42}{2}}{2} = \frac{41+21}{2} = \boxed{\textbf{(B)}~31}$

 


	\item (\textit{AMC 8 2015}) A triangle with vertices as $A=(1,3)$, $B=(5,1)$, and $C=(4,4)$ is plotted on a $6\times5$ grid. What fraction of the grid is covered by the triangle?


 
                    \begin{figure}[H]
                        \centering
                        \includegraphics[width=0.5\textwidth]{Problem/AMC_8_AMC_2015_19_Problem_Figure_1.png}
                    \end{figure}
                    


 $\textbf{(A) }\frac{1}{6} \qquad \textbf{(B) }\frac{1}{5} \qquad \textbf{(C) }\frac{1}{4} \qquad \textbf{(D) }\frac{1}{3} \qquad \textbf{(E) }\frac{1}{2}$


 
 \textbf{Answer: }A



\textbf{Solution 1}

The area of $\triangle ABC$ is equal to half the product of its base and height.  By the Pythagorean Theorem, we find its height is $\sqrt{1^2+2^2}=\sqrt{5}$, and its base is $\sqrt{2^2+4^2}=\sqrt{20}$.  We multiply these and divide by $2$ to find the area of the triangle is $\frac{\sqrt{5 \cdot 20}}2=\frac{\sqrt{100}}2=\frac{10}2=5$.  Since the grid has an area of $30$, the fraction of the grid covered by the triangle is $\frac 5{30}=\boxed{\textbf{(A) }\frac{1}{6}}$.

 


 



\textbf{Solution 2}

Note angle $\angle ACB$ is right; thus, the area is $\sqrt{1^2+3^2} \times \sqrt{1^2+3^2}\times \dfrac{1}{2}=10 \times \dfrac{1}{2}=5$; thus, the fraction of the total is $\dfrac{5}{30}=\boxed{\textbf{(A)}~\dfrac{1}{6}}$.

 


 



\textbf{Solution 3 (Shoelace Theorem)}

By the \href{https://artofproblemsolving.com/wiki/index.php?title=Shoelace\_Theorem}{Shoelace Theorem}, the area of $\triangle ABC=|\dfrac{1}{2}(15+4+4-1-20-12)|=|\dfrac{1}{2}(-10)|=5$.

 This means the fraction of the total area is $\dfrac{5}{30}=\boxed{\textbf{(A)}~\dfrac{1}{6}}$.

 



\textbf{Solution 4 (easiest solution)}

The smallest rectangle that follows the grid lines and completely encloses $\triangle ABC$ has an area of $12$, where $\triangle ABC$ splits the rectangle into four triangles. The area of $\triangle ABC$ is therefore $12 - (\frac{4 \cdot 2}{2}+\frac{3 \cdot 1}{2}+\frac{3 \cdot 1}{2}) = 12 - (4 + \frac{3}{2} + \frac{3}{2}) = 12 - 7 = 5$. That means that $\triangle ABC$ takes up $\frac{5}{30} = \boxed{\textbf{(A)}~\frac{1}{6}}$ of the grid.

 



\textbf{Solution 5 (Pick's Theorem, very much recommended to learn this)}

Using \href{https://artofproblemsolving.com/wiki/index.php?title=Pick\%27s\_Theorem}{Pick's Theorem}, the area of the triangle is $4 + \dfrac{4}{2} - 1=5$. Therefore, the triangle takes up $\dfrac{5}{30}=\boxed{\textbf{(A)}~\frac{1}{6}}$ of the grid.

 


 



\textbf{Solution 6 (Heron's Formula, Not Recommended)}

We can find the lengths of the sides by using the \href{https://artofproblemsolving.com/wiki/index.php?title=Pythagorean\_Theorem}{Pythagorean Theorem}. Then, we apply \href{https://artofproblemsolving.com/wiki/index.php?title=Heron\%27s\_Formula}{Heron's Formula} to find the area. 
 \[\sqrt{\left(\frac{\sqrt{10}+\sqrt{10}+2\sqrt{5}}{2}\right)\left(\frac{\sqrt{10}+\sqrt{10}+2\sqrt{5}}{2}-\sqrt{10}\right)\left(\frac{\sqrt{10}+\sqrt{10}+2\sqrt{5}}{2}-\sqrt{10}\right)\left(\frac{\sqrt{10}+\sqrt{10}+2\sqrt{5}}{2}-2\sqrt{5}\right)}.\]This simplifies to
 \[\sqrt{\left(\sqrt{10}+\sqrt{5}\right)\left(\sqrt{10}+\sqrt{5}-\sqrt{10}\right)\left(\sqrt{10}+\sqrt{5}-\sqrt{10}\right)\left(\sqrt{10}+\sqrt{5}-2\sqrt{5}\right)}.\]Again, we simplify to get
 \[\sqrt{\left(\sqrt{10}+\sqrt{5}\right)\left(\sqrt{5}\right)\left(\sqrt{5}\right)\left(\sqrt{10}-\sqrt{5}\right)}.\]The middle two terms inside the square root multiply to $5$, and the first and last terms inside the square root multiply to $\sqrt{10}^2-\sqrt{5}^2=10-5=5.$ This means that the area of the triangle is
 \[\sqrt{5\cdot 5}=5.\]The area of the grid is $6\cdot 5=30.$ Thus, the answer is $\frac{5}{30}=\boxed{\textbf{(A) }\frac{1}{6}}$.

 


 



\textbf{Solution 7 (Simple Deduction)}

First, count the number of 1x1 squares inside the main triangle, both fully enclosed and partially (you should count 10). Then, upon closer inspection, most of the squares that are not a single unit on the triangle can be created by connecting another shape. The only exception is one square that is a single unit and one that would need 2 shapes connected to it to make a single unit. On average, you need to connect 2 squares to make a unit. Knowing this, if there are 10 squares and you require 2 shapes to make a unit, 10 divided by 2 equals 5, which is the area.5 is 1/6 of 30(the total of the graph) and so the final answer is $\frac{5}{30}=\boxed{\textbf{(A) }\frac{1}{6}}$.

 -Themathnerd3.14

 


	\item (\textit{AMC 8 2015}) Ralph went to the store and bought $12$ pairs of socks for a total of $\$24$. Some of the socks he bought cost $\$1$ a pair, some of the socks he bought cost $\$3$ a pair, and some of the socks he bought cost $\$4$ a pair. If he bought at least one pair of each type, how many pairs of $\$1$ socks did Ralph buy?


 $\textbf{(A) } 4 \qquad \textbf{(B) } 5 \qquad \textbf{(C) } 6 \qquad \textbf{(D) } 7 \qquad \textbf{(E) } 8$


 
 \textbf{Answer: }D



\textbf{Solution 1}

So, let there be $x$ pairs of $\$1$ socks, $y$ pairs of $\$3$ socks, and $z$ pairs of $\$4$ socks.

 We have $x+y+z=12$, $x+3y+4z=24$, and $x,y,z \ge 1$.

 Now, we subtract to find $2y+3z=12$, and $y,z \ge 1$.It follows that $2y$ is a multiple of $3$ and $3z$ is a multiple of $3$. Since sum of 2 multiples of 3 = multiple of 3, so we must have $2y=6$.

 Therefore, $y=3$, and it follows that $z=2$. Now, $x=12-y-z=12-3-2=\boxed{\textbf{(D)}~7}$, as desired.

 



\textbf{Solution 2}

Since the total cost of the socks was $\$24$ and Ralph bought $12$ pairs, the average cost of each pair of socks is $\frac{\$24}{12} = \$2$.

 There are two ways to make packages of socks that average to $\$2$. You can have:

 $\bullet$ Two $\$1$ pairs and one $\$4$ pair (package adds up to $\$6$)

 $\bullet$ One $\$1$ pair and one $\$3$ pair (package adds up to $\$4$)

 Now, we need to solve
 \[6a+4b=24,\]where $a$ is the number of $\$6$ packages and $b$ is the number of $\$4$ packages. We see our only solution (that has at least one of each pair of sock) is $a=2, b=3$, which yields the answer of $2\times2+3\times1 = \boxed{\textbf{(D)}~7}$.

 



\textbf{Solution 3}

Since there are 12 pairs of socks, and Ralph bought at least one pair of each, there are $12-3=9$ pairs of socks left. Also, the sum of the three pairs of socks is $1+3+4=8$. This means that there are $24-8=16$ dollars left. If there are only $1$ dollar socks left, then we would have $9\cdot1=9$ dollars wasted, which leaves $7$ more dollars. If we replace one pair with a $3$ dollar pair, then we would waste an additional $2$ dollars. If we replace one pair with a $4$ dollar pair, then we would waste an additional $3$ dollars. The only way $7$ can be represented as a sum of $2$s and $3$s is $2+2+3$. If we change $3$ pairs, we would have $6$ pairs left. Adding the one pair from previously, we have $\boxed{(\text{D})~7}$ pairs.

 



\textbf{Solution 4}

Let the amount of $1$ dollar socks be $a$, $3$ dollar socks be $b$, and $4$ dollar socks be $c$. We then know that $a+b+c=12$ and $a+3b+4c=24$. We can make $a+b+c=12$ into $a=12-b-c$ and then plug that into the other equation, producing $12-b-c+3b+4c=24$ which simplifies to $2b+3c=12$. It's not hard to see $b=3$ and $c=2$. Now that we know $b$ and $c$, we know that $a=7$, meaning the number of $1$ dollar socks Ralph bought is $\boxed{\textbf{(D)} 7}$.

 



\textbf{Solution 5 (Guess and check)}

If Ralph bought one sock of each kind, he already used $\$8$, so there are $\$16$ left and 9 socks. If we split the $\$16$ into four $\$4$ sections, (as it is the smallest possible number that 1, 3, 4, can make in different ways that in all use at least each of the numbers once,) if Ralph bought a $\$3$ pair, he would need to buy a $\$1$ pair in order for it to add up to a multiple of four. Similarly, if Ralph bought a $\$1$ pair, he would either need to buy three $\$1$ pairs or a $\$3$ pair. If Ralph bought a $\$4$ pair, it would already make a group. Now, the problem is just how we can split 9 into 4 groups of 1, 2, or 4. We clearly see that $1 + 2 + 2 + 4 = 9$, or a $\$4$ pair, two $\$3$ pairs, and six $\$1$ pairs. Because we subtracted the necessary one of each kind, there are two $\$4$ pairs, three $\$3$ pairs, and seven $\$1$ pairs. Therefore, the number of $\$1$ pairs Ralph bought is $\boxed{\textbf{(D)}~7}$.~strongstephen

 



\textbf{Solution 6}

To simplify the problem, we can reduce the price of each sock by $\$1$, resulting in the prices being $\$0$, $\$2$ and $\$3$. This means the total price Ralph spends becomes $\$24-12\times\$1=\$12$. Now the $\$0$ socks aren't worth anything, but we still need to keep in mind there are at least one of them, so the question becomes: at most 11 of $\$2$ and $\$3$ socks add up to $\$12$. It is obvious that $3\times\$2+2\times\$3=\$12$, so the number of $\$0$ socks (actually $\$1$ socks) is $12-3-2=7$.

 Note: this is very similar to Solution 1, but a non algebraic way of thinking about it.

 -bbidev

 


	\item (\textit{AMC 8 2015}) In the given figure hexagon $ABCDEF$ is equiangular, $ABJI$ and $FEHG$ are squares with areas $18$ and $32$ respectively, $\triangle JBK$ is equilateral and $FE=BC$. What is the area of $\triangle KBC$?


 
                    \begin{figure}[H]
                        \centering
                        \includegraphics[width=0.5\textwidth]{Problem/AMC_8_AMC_2015_21_Problem_Figure_1.png}
                    \end{figure}
                    


 $\textbf{(A) }6\sqrt{2}\qquad\textbf{(B) }9\qquad\textbf{(C) }12\qquad\textbf{(D) }9\sqrt{2}\qquad\textbf{(E) }32$


 
 \textbf{Answer: }C



\textbf{Solution 1}

Clearly, since $\overline{FE}$ is a side of a square with area $32$, $\overline{FE} = \sqrt{32} = 4 \sqrt{2}$. Now, since $\overline{FE} = \overline{BC}$, we have $\overline{BC} = 4 \sqrt{2}$.

 We know that $\overline{JB}$ is a side of a square with area $18$, so $\overline{JB} = \sqrt{18} = 3 \sqrt{2}$. Since $\Delta JBK$ is equilateral, $\overline{BK} = \overline{JB} = 3 \sqrt{2}$.

 Lastly, $\Delta KBC$ is a right triangle. We see that $\angle JBA + \angle ABC + \angle CBK + \angle KBJ = 360^{\circ} \rightarrow 90^{\circ} + 120^{\circ} + \angle CBK + 60^{\circ} = 360^{\circ} \rightarrow \angle CBK = 90^{\circ}$, so $\Delta KBC$ is a right triangle with legs $3 \sqrt{2}$ and $4 \sqrt{2}$. Now, its area is $\dfrac{3 \sqrt{2} \cdot 4 \sqrt{2}}{2} = \dfrac{24}{2} = \boxed{\textbf{(C)}~12}$.

 



\textbf{Solution 2}

Since $\overline{FE} = \sqrt{32}$, and $\overline{FE} = \overline{BC}$, $\overline{BC} = 4\sqrt{2}$. Meanwhile, $\overline{JB} = 3\sqrt{2}$, and since $\triangle JBK$ is equilateral, $\overline{BK} = 3\sqrt{2}$. If $ABCDEF$ is equiangular, $\angle ABC = \frac{180 \cdot (n-2)}{n} = 120^{\circ}$, where $n$ is the number of sides of the shape. Adding all the angles around $B$ gives $270^{\circ}$, so $\angle KBC = 360 - 90 = 270^{\circ}$. Because $\triangle KBC$ is right, the area of $\triangle KBC = \frac{4\sqrt{2} \cdot 3\sqrt{2}} {2} = 12$. Therefore, the answer is $\boxed{\textbf{(C)}~12}$. ~strongstephen

 Note: Another, more complicated way to determine if angle KBC is right is to extend line AB to a point L on triangle KBC. Now, since we know that angle IJB is right, and we know that line IJ now acts as a transversal to lines IJ and AL. That means angle IJB (90 degrees) must be equal to angle JBL by the alternate exterior angle rule. Since we know that triangle JKB is equilateral, angle JBK will be 60 degrees, hence angle KBL will be 90 - 60 = 30 degrees. Now, we move on to figure out ABCDEF. Through the solution above, we know that angle ABC is 120 degrees, so we can use the same line AL, which will result in the formation of a linear pair of angles: angles ABC and CBL. Since the angles in a linear pair always add up to 180, and angle ABC is 120, angle CBL = 180 - 120 = 60 degrees. Notice that angles CBL and KBL together make up angle KBC. Hence, by the value of 30 degrees derived for angle KBL, angle KBC = angle KBL + angle CBL = 30 degrees + 60 degrees = 90 degrees.

 



\textbf{Solution 3}

Square ABJI's side length is $\sqrt{18}=3\sqrt{2}$, and since Triangle JBK is equilateral, that means that the side length of Triangle JBK would be $3\sqrt{2}$, which means that KB, the height of Triangle KBC would be $3\sqrt{2}$. Square FEHG's side length would then be $\sqrt{32}=4\sqrt{2}$, and because line FE = line BC, The base of Triangle KBC would be $4\sqrt{2}$. The formula to finding the area of a triangle is the $\frac{Height \cdot Base}{2}$, which means that Triangle KBC's area would just be $\frac{3\sqrt{2}\cdot4\sqrt{2}}{2}$, and the answer is $\boxed{12}.$

 


 


	\item (\textit{AMC 8 2015}) On June $1$, a group of students are standing in rows, with $15$ students in each row. On June $2$, the same group is standing with all of the students in one long row. On June $3$, the same group is standing with just one student in each row. On June $4$, the same group is standing with $6$ students in each row. This process continues through June $12$ with a different number of students per row each day. However, on June $13$, they cannot find a new way of organizing the students. What is the smallest possible number of students in the group?


 $\textbf{(A) } 21 \qquad \textbf{(B) } 30 \qquad \textbf{(C) } 60 \qquad \textbf{(D) } 90 \qquad \textbf{(E) } 1080$


 
 \textbf{Answer: }C



\textbf{Solution 1}

Since we know the number must be a multiple of $15$, we can eliminate $A$. We also know that after $12$ days, the students can't find any more arrangements, meaning the number has $12$ factors. Now, we just list the factors of every number, starting with $30$:
 \[30=1\cdot30, 2\cdot15, 3\cdot10, 5\cdot6\]
 \[60=1\cdot60, 2\cdot30, 3\cdot20, 4\cdot15, 5\cdot12, 6\cdot10\]$60$ has $12$ factors, so the answer is $\boxed{\textbf{(C) } 60}$.

 Note - A faster way to find the factors of a number is really helpful in competitions, as you don't want to waste time listing them all.What we can do is find the prime factorization of a number. We'll use $60$ as an example here, so we take $60 = 2 \cdot 2 \cdot 3 \cdot 5$ or $60 = 2^2 \cdot 3^1 \cdot 5^1$. Now, we just add one to the power of each factor. So since we have powers of $2$, $1$, and $1$, we turn those into $3$, $2$, and $2$. Finally, we just multiply them together to get $3 \cdot 2 \cdot 2 = 12$ as the amount of factors in 60.

 



\textbf{Solution 2}

We know that the number of students in the group has to be a multiple of 15, so we can eliminate A. However, on June 13, they cannot find a new way of organizing the students meaning that the number has only 12 factors.The prime factorization of 60 is equal to $2^2 \times 3^1 \times 5^1$Using the factor counting formula, we have the add one to each exponent and multiply them together to find the number of factors.$(2+1) \times (1+1) \times (1+1) = 3 \times 2 \times 2 = 12$ factors making $\boxed{\textbf{(C) } 60}$ the answer.

 



\textbf{Solution 3 (Combination of Solutions)}

On the 1st day, all the students are in an arrangement such that there are $15$ students per row. On the 4th day, all the students are in an arrangement such that there are $6$ students per row. This tells us that the \# of students is a multiple of both $15$ and $6$. Since we can arrange the students for $12$ days, this tells us that the \# of students has $12$ factors. We can also tell that the \# of students is a multiple of $30$, as $30$ is the $LCM$ of $6$ and $15$. 

 
We now test answer choice B.

 
$30$ = $2*3*5$  =>  30 has $2*2*2 = 8$ factors.

 
Let's test answer choice C.

 
$60$ = $2^2*3*5$  =>  60 has $3*2*2 = 12$ factors.

 
We know that $\boxed{\textbf{(C) } 60}$ is our answer.

 


	\item (\textit{AMC 8 2015}) Tom has twelve slips of paper which he wants to put into five cups labeled $A$, $B$, $C$, $D$, $E$. He wants the sum of the numbers on the slips in each cup to be an integer. Furthermore, he wants the five integers to be consecutive and increasing from $A$ to $E$. The numbers on the papers are $2, 2, 2, 2.5, 2.5, 3, 3, 3, 3, 3.5, 4,$ and $4.5$. If a slip with $2$ goes into cup $E$ and a slip with $3$ goes into cup $B$, then the slip with $3.5$ must go into what cup?


 $\textbf{(A) } A \qquad \textbf{(B) } B \qquad \textbf{(C) } C \qquad \textbf{(D) } D \qquad \textbf{(E) } E$


 
 \textbf{Answer: }D



\textbf{Solution}

The numbers have a sum of $6+5+12+4+8=35$, which averages to $7$, which means $A, B, C, D,$ and$E$ will have the values $5, 6, 7, 8$ and $9$, respectively. Now, it's the process of elimination: Cup $A$will have a sum of $5$, so putting a $3.5$ slip in the cup will leave $5-3.5=1.5$; however,all of our slips are bigger than $1.5$, so this is impossible. Cup $B$ has a sum of $6$, but we are told that it already has a $3$ slip, leaving $6-3=3$, which is too small for the$3.5$ slip. Cup $C$ is a little bit trickier but still manageable. It must have a value of $7$, so adding the $3.5$ slip leaves room for $7-3.5=3.5$. This looks good at first as we have slips smaller than that, but upon closer inspection, we see that no slip fits exactly. And the smallest sum of two slips is $2+2=4$, which is too big, so this case is also impossible. Cup $E$ has a sum of $9$, but we are told it already has a $2$ slip. So, weare left with $9-2=7$, which is identical to Cup C and thus also impossible.With all other choices removed, we are left with the answer: Cup $\boxed{\textbf{(D)}~D}$.

 


	\item (\textit{AMC 8 2015}) A baseball league consists of two four-team divisions. Each team plays every other team in its division $N$ games. Each team plays every team in the other division $M$ games with $N>2M$ and $M>4$. Each team plays a $76$-game schedule. How many games does a team play within its own division?


 $\textbf{(A) } 36 \qquad \textbf{(B) } 48 \qquad \textbf{(C) } 54 \qquad \textbf{(D) } 60 \qquad \textbf{(E) } 72$


 
 \textbf{Answer: }B



\textbf{Solution 1}

On one team they play $3N$ games in their division and $4M$ games in the other.  This gives $3N+4M=76$.

 Since $M>4$ we start by trying M=5. This doesn't work because $56$ is not divisible by $3$.

 Next, $M=6$ does not work because $52$ is not divisible by $3$.

 We try $M=7$ does work by giving $N=16$ ,$~M=7$ and thus $3\times 16=\boxed{\textbf{(B)}~48}$ games in their division.

 $M=10$ seems to work, until we realize this gives $N=12$, but $N>2M$ so this will not work.

 



\textbf{Solution 2}

$76=3N+4M > 10M$, giving $M \le 7$.Since $M>4$, we have $M=5,6,7$.Since $4M$ is $1$ $\pmod{3}$, we must have $M$ equal to $1$ $\pmod{3}$, so $M=7$.

 This gives $3N=48$, as desired. The answer is $\boxed{\textbf{(B)}~48}$.

 



\textbf{Solution 3}

Notice that each team plays $N$ games against each of the three other teams in its division. So that's $3N$.

 Since each team plays $M$ games against each of the four other teams in the other division, that's $4M$.

 So $3N+4M=76$, with $M>4, N>2M$. 

 Let's start by solving this Diophantine equation. In other words, $N=\frac{76-4M}{3}$. 

 So $76-4M\equiv0 \pmod{3}$ (remember: $M$ must be divisible by 3 for $N$ to be an integer!). Therefore, after reducing $76$ to $1$ and $-4M$ to $2M$ (we are doing things in $\pmod{3}$), we find that $M\equiv1 \pmod{3}$. 

 Since $M>4$, so the minimum possible value of $M$ is $7$. However, remember that $N>2M$! To find the greatest possible value of M, we assume that $N=2M$ and that is the upper limit of $M$ (excluding that value because $N>2M$). Plugging $N=2M$ in, $10M=76$. So $M<7.6$. Since you can't have $7.6$ games, we know that we can only check $M=7$ since we know that since $M>4, M<7.6, M\equiv1 \pmod{3}$. After checking $M=7$, we find that it will work.

 So $M=7, N=16$. So each team plays 16 games against each team in its division. Since they are asking for games in it division, which equals $3n = 48$. Select $\boxed{\textbf{(B)}~48}$. (P.S: This might be too complicated).

 


	\item (\textit{AMC 8 2015}) One-inch squares are cut from the corners of this $5$ inch square. What is the area in square inches of the largest square that can fit into the remaining space?


 
                    \begin{figure}[H]
                        \centering
                        \includegraphics[width=0.5\textwidth]{Problem/AMC_8_AMC_2015_25_Problem_Figure_1.png}
                    \end{figure}
                    


 $\textbf{(A) } 9\qquad \textbf{(B) } 12\frac{1}{2}\qquad \textbf{(C) } 15\qquad \textbf{(D) } 15\frac{1}{2}\qquad \textbf{(E) } 17$


 
 \textbf{Answer: }C



\textbf{Solution 1}

We can draw a diagram as shown:
                    \begin{figure}[H]
                        \centering
                        \includegraphics[width=0.5\textwidth]{Solution/AMC_8_AMC_2015_25_Solution_Figure_1.png}
                    \end{figure}
                     Let us focus on the $4$ big triangles that, together with the inscribed square, fill in the large square. If we zoom in on one of the four triangles, we can see that it is composed of a small unit square in the corner of the large square and two triangles, one smaller than the other. We are going to focus specifically on the smaller of the two triangles. This triangle is similar to the big triangle itself by $\mathrm{AA}$ similarity (because the two sides of a square are parallel. To prove this fact, draw a diagonal of the square and find congruent triangles). Let the shorter leg of the big triangle be $x$; then $\dfrac{x}{x-1}=\dfrac{5-x}{1}$.
 \[x=-x^2+6x-5\]
 \[x^2-5x+5=0\]
 \[x=\dfrac{5\pm \sqrt{(-5)^2-(4)(1)(5)}}{2}\]
 \[x=\dfrac{5\pm \sqrt{5}}{2}\]Thus, $x=\dfrac{5-\sqrt{5}}{2}$, because by symmetry, $x < \dfrac52$. Note that the other solution we got, namely, $x=\dfrac{5+\sqrt 5} 2$, is the length of the segment $5-x$. $x$ and $5-x$ together sum to $5$, the side of the length of the large square, and similarly, the sum of the solutions is $5$. This solution is a result of the symmetry of the problem; if we had set the longer leg of the big triangle to be $x$, then we would solve the same quadratic to find the same roots, the only difference being that we take the other root.

 This means the area of each triangle is $\dfrac{5-\sqrt{5}}{2}*(5-\dfrac{5-\sqrt{5}}{2})*\dfrac{1}{2}=\dfrac{5}{2}$.Thus, the area of the square is $25-(4*\dfrac{5}{2})=15\implies \boxed{\textbf{(C)}}$.

 Don't use this in a contest scenario, and only use it when practicing math skills. :)

 


 



\textbf{Solution 2 (Contest Solution, most recommended and fastest)}

We draw a square as shown:

 
                    \begin{figure}[H]
                        \centering
                        \includegraphics[width=0.5\textwidth]{Solution/AMC_8_AMC_2015_25_Solution_Figure_2.png}
                    \end{figure}
                    

 
We want to find the area of the biggest square. The area of this square is composed of the center white square and the four red triangles. Because the inner square has an area of $(5-2)$, $3$, squared, $9$, it also has a length of $\sqrt{9}=3$. The heights of each of the red triangles are 1 (because the gray squares have lengths of one), and the area of one triangle is namely $\frac{3 \cdot 1}{2}$. Thus, the combined area of the four triangles is $4 \cdot \frac 32=6$. Furthermore, the area of the smaller square is $9$. We add these to see that the area of the large square is $9+6=15\implies \boxed{\textbf{(C)}}$.

 



\textbf{Solution 3}

Let us find the area of the triangles and the unit squares: on each side, there are two triangles. They both have one leg of length $1$, and let's label the other legs $x$ for one of the triangles and $y$ for the other. Note that $x + y = 3$.The area of each of the triangles is $\frac{x}{2}$ and $\frac{y}{2}$, and there are $4$ of each. So now, we need to find $4\left(\frac{x}{2}\right) + 4\left(\frac{y}{2}\right)$. Expanding out the fractions, we get $4\left(\frac{x}{2}\right) + 4\left(\frac{y}{2}\right) = 2x + 2y = 2(x + y) = 2(3) = 6$.So, the area of the square we need is $25- (4+6) = 15\implies \boxed{\textbf{(C)}}$.

 



\textbf{Solution 4}

Proceed as in solution 1 to get $x^2 - 5x + 5 = 0$.

 Notice that the side length of the square is $\sqrt{x^2+(5-x)^2}$ by the Pythagorean Theorem. Thus, the area of the square is
 \[\left(\sqrt{x^2+(5-x)^2}\right)^2\]
 \[= x^2 + x^2 - 10x + 25\]
 \[= (2x^2 - 10x + 10) + 15\]
 \[= 2(x^2 - 5x + 5) + 15\]
 \[= 2(0) + 15\]
 \[= \boxed{15}\]
 \[\implies \boxed{\textbf{(C)}}\]

 



\textbf{Solution 5 (Fast with no explanations)}

By inspection, we can see that the area of the largest square would be the area of the middle square formed by the vertices of the shaded squares plus half the combined area of the four smaller rectangles. Our desired answer is $9 + \frac{3\cdot4}{2} = 15\implies \boxed{\textbf{(C)}}$.

 


\end{enumerate}